% !PRINTED 516
\weberpage{516}

Denn wenn $\alpha$ und $\beta$ beide nicht durch $\pi$ theilbar, also beide relativ prim zu $\pi$ sind, so ist nach \S.\,139, 4. auch $\alpha \beta$ relativ prim zu $\pi$.

Es ergiebt sich daraus durch wiederholte Anwendung, dass, wenn ein Product aus \gesperrt{mehreren} Factoren durch $\pi$ theilbar ist, mindestens einer der Factoren durch $\pi$ theilbar sein muss.

3. \gesperrt{Die kleinste natürliche ganze Zahl $p$, die durch ein Primfunctional $\pi$ in $\Omega$ theilbar ist, ist eine natürliche Primzahl, und die absolute Norm von $\pi$ ist eine Potenz von $p$. Jede durch $\pi$ theilbare ganze rationale Zahl ist auch durch $p$ theilbar.}

Nach \S.\,138, 13. giebt es natürliche Zahlen, die durch $\pi$ theilbar sind. Die kleinste unter ihnen, $p$, kann nicht in zwei natürliche Factoren, die grösser als 1 sind, zerlegbar sein; denn ist $p = p_1 p_2$, so muss entweder $p_1$ oder $p_2$ durch $\pi$ theilbar sein (nach 2.). Ist aber keiner der Factoren $p_1, p_2 = 1$, so sind sie beide kleiner als $p$, was der Voraussetzung über $p$ widerspricht. Folglich ist $p$ eine natürliche Primzahl. Dass jede durch $\pi$ theilbare natürliche Zahl $m$ durch $p$ theilbar ist, haben wir schon im \S.\,138, 13. bewiesen.

Setzen wir nun

\weq{1}{p = \pi \omega,}

so ist $\omega$ ein ganzes Functional, und wenn wir beiderseits die absoluten Normen nehmen, so folgt, da die absolute Norm einer natürlichen Zahl die $n^{\text{te}}$ Potenz dieser Zahl ist:

\weq{2}{p^n = N_a (\pi) N_a (\omega).}

Hieraus folgt der zweite Theil unseres Satzes, dass die natürliche ganze Zahl $N_a (\pi)$ eine Potenz von $p$ ist.

Setzen wir demnach

\weq{3}{N_a (\pi) = p^f,}

so ist $f$ eine Zahl, die nur einen der Werthe $1, 2, 3, \dots, n$ haben kann, wenn $n$ den Grad des Körpers $\Omega$ bedeutet.

Die Zahl $f$ heisst der \gesperrt{Grad} des Primfunctionals $\pi$.

\begin{center}
\S.\,141.

\gesperrt{Zerlegung der ganzen und gebrochenen Functionale in Primfactoren.}
\end{center}

1. \gesperrt{Jedes von Null verschiedene ganze Functional $\omega$ des Körpers $\Omega$, das keine Einheit ist, ist durch ein Primfunctional theilbar.}
